% ============================================================
%  PREÀMBUL — Dossier 1: Circuits amb resistències (4t ESO)
% ============================================================
\usepackage[utf8]{inputenc}
\usepackage[T1]{fontenc}
\usepackage{textcomp}
% El punt volat (·) es dibuixa amb el símbol matemàtic \cdot: així es veu igual
% a qualsevol tipus de lletra i amb qualsevol motor (pdfLaTeX, XeTeX/Tectonic, LuaTeX).
\usepackage{newunicodechar}
\newunicodechar{·}{\ensuremath{\cdot}}
\newcommand{\sep}{\ \ensuremath{\cdot}\ }
\usepackage[catalan]{babel}
\usepackage[sc]{mathpazo}
\usepackage[a4paper,left=2cm,right=2cm,top=2.3cm,bottom=2.2cm]{geometry}
\usepackage{microtype}
\usepackage{xcolor}
\usepackage{amsmath,amssymb}
\usepackage{siunitx}
\sisetup{output-decimal-marker={,},group-minimum-digits=5,group-separator={\,},
         detect-all=false,per-mode=symbol}
% Amb XeTeX/LuaTeX (p. ex. Tectonic), el símbol Ω que dona siunitx no existeix a la lletra:
% el dibuixem amb la lletra grega de les matemàtiques.
\usepackage{iftex}
\ifPDFTeX\else\DeclareSIUnit{\ohm}{\ensuremath{\mathrm{\Omega}}}\fi
\usepackage[european]{circuitikz}
\usepackage[most]{tcolorbox}
\usepackage{booktabs,array,enumitem,fancyhdr,titlesec,needspace,pifont,graphicx}
\usepackage[hidelinks]{hyperref}

\linespread{1.08}
\setlength{\parindent}{0pt}
\setlength{\parskip}{0.5em}
\setlist{nosep,leftmargin=1.4em}

% ---------- Colors ----------
\definecolor{blau}{HTML}{1F5F99}
\definecolor{verd}{HTML}{2E7D32}
\definecolor{taronja}{HTML}{E07B00}
\definecolor{vermell}{HTML}{B3261E}
\definecolor{lila}{HTML}{6A3FA0}

% ---------- Títols ----------
\titleformat{\section}{\Large\bfseries\color{blau}}{Bloc \thesection}{0.7em}{}[\vspace{-0.4em}\color{blau}\titlerule]
\titleformat{\subsection}{\large\bfseries\color{blau}}{}{0em}{}
\titlespacing*{\section}{0pt}{0pt}{1em}
\titlespacing*{\subsection}{0pt}{1.2em}{0.4em}

% ---------- Capçalera i peu ----------
\pagestyle{fancy}
\fancyhf{}
\renewcommand{\headrulewidth}{0.3pt}
\fancyhead[L]{\small\textsc{Electrònica analògica}\sep 4t ESO}
\fancyhead[R]{\small Dossier 1\sep Circuits amb resistències\ifsolucions\sep\textbf{SOLUCIONARI}\fi}
\fancyfoot[C]{\small\thepage}

% ---------- Unitats i símbols de circuit ----------
\ctikzset{bipoles/length=1.25cm, bipoles/resistor/height=0.28, bipoles/resistor/width=0.7}

% ---------- Quadres de colors ----------
\tcbset{boxrule=0.6pt, arc=3pt, left=7pt, right=7pt, top=5pt, bottom=5pt,
        fonttitle=\bfseries, breakable, before skip=0.8em, after skip=0.8em}
\newtcolorbox{idea}{colback=blau!6, colframe=blau, title={La idea clau}}
\newtcolorbox{metode}[1]{colback=lila!6, colframe=lila, title={Mètode: #1}}
\newtcolorbox{exemple}[1]{colback=verd!7, colframe=verd, title={Exemple resolt: #1}}
\newtcolorbox{guiada}[1]{colback=taronja!8, colframe=taronja, title={Practiquem junts: #1}}
\newtcolorbox{atencio}{colback=vermell!6, colframe=vermell, title={Compte amb l'error típic}}
\newtcolorbox{comprova}{colback=yellow!12, colframe=yellow!55!black, title={Comprova-ho}}

% ---------- Nivells d'exercici ----------
\newcommand{\estrella}{\ding{72}}
\newcommand{\nivell}[1]{%
  \ifcase#1\or{\color{verd}\estrella}\or{\color{taronja}\estrella\estrella}\or{\color{vermell}\estrella\estrella\estrella}\fi}

\newcounter{exercici}[section]
\renewcommand{\theexercici}{\thesection.\arabic{exercici}}
\newenvironment{exercici}[2][8]{%
  \par\bigskip\needspace{#1\baselineskip}\noindent
  \refstepcounter{exercici}\textbf{\theexercici}\enspace\nivell{#2}\enspace\ignorespaces}%
  {\par}

% ---------- Solucions (espai en blanc o quadre de solució) ----------
%  \resposta{espai en blanc}{solució}
\newcommand{\resposta}[2]{%
  \ifsolucions
    \par\smallskip
    \begin{tcolorbox}[colback=verd!6, colframe=verd!60!black, boxrule=0.5pt, arc=2pt,
                      left=5pt, right=5pt, top=3pt, bottom=3pt, before skip=0.3em, after skip=0.6em]
      {\small\textbf{Solució.}\ #2}
    \end{tcolorbox}
  \else
    \vspace*{#1}
  \fi}

% ---------- Passos del mètode (per a les pràctiques guiades) ----------
\newcommand{\pas}[3]{\par\noindent\textbf{\textcolor{taronja}{\ding{\numexpr171+#1\relax}~#2}}\quad #3\par\smallskip}

% ---------- Taula de magnituds ----------
\newenvironment{taulamag}{%
  \par\smallskip\renewcommand{\arraystretch}{1.7}\centering
  \begin{tabular}{@{}l *{4}{>{\centering\arraybackslash}p{2.5cm}}@{}}
  \toprule
  \textbf{Element} & $V$ (\si{\volt}) & $I$ (\si{\ampere}) & $R$ (\si{\ohm}) & $P$ (\si{\watt})\\
  \midrule}%
  {\bottomrule\end{tabular}\par\smallskip}

% ---------- Dibuixos de circuits ----------
\newcommand{\simbol}[1]{\begin{circuitikz}[scale=0.8,transform shape]\draw (0,0) to[#1] (1.6,0);\end{circuitikz}}

% Una resistència: V, R
\newcommand{\csimple}[2]{%
\begin{circuitikz}[scale=0.9,transform shape]
 \draw (0,0) to[battery1,l_={#1}] (0,2.2) to[R,l={#2}] (3.5,2.2) to[short] (3.5,0) to[short] (0,0);
\end{circuitikz}}

% Dues resistències en sèrie
\newcommand{\cserie}[3]{%
\begin{circuitikz}[scale=0.9,transform shape]
 \draw (0,0) to[battery1,l_={#1}] (0,2.2) to[R,l={#2}] (3.2,2.2) to[R,l={#3}] (6.4,2.2)
       to[short] (6.4,0) to[short] (0,0);
\end{circuitikz}}

% Tres resistències en sèrie
\newcommand{\cserietres}[4]{%
\begin{circuitikz}[scale=0.8,transform shape]
 \draw (0,0) to[battery1,l_={#1}] (0,2.2) to[R,l={#2}] (3,2.2) to[R,l={#3}] (6,2.2) to[R,l={#4}] (9,2.2)
       to[short] (9,0) to[short] (0,0);
\end{circuitikz}}

% Dues resistències en paral·lel
\newcommand{\cparalel}[3]{%
\begin{circuitikz}[scale=0.9,transform shape]
 \draw (0,0) to[battery1,l_={#1}] (0,2.6) to[short] (3,2.6) to[R,l={#2}] (3,0) to[short] (0,0);
 \draw (3,2.6) to[short,*-] (6,2.6) to[R,l={#3}] (6,0) to[short,-*] (3,0);
\end{circuitikz}}

% Tres resistències en paral·lel
\newcommand{\cparaleltres}[4]{%
\begin{circuitikz}[scale=0.85,transform shape]
 \draw (0,0) to[battery1,l_={#1}] (0,2.6) to[short] (3,2.6) to[R,l={#2}] (3,0) to[short] (0,0);
 \draw (3,2.6) to[short,*-] (6,2.6) to[R,l={#3}] (6,0) to[short,-*] (3,0);
 \draw (6,2.6) to[short,*-] (9,2.6) to[R,l={#4}] (9,0) to[short,-*] (6,0);
\end{circuitikz}}

% Mixt A: R1 en sèrie amb (R2 en paral·lel amb R3)
\newcommand{\cmixa}[4]{%
\begin{circuitikz}[scale=0.85,transform shape]
 \draw (0,0) to[battery1,l_={#1}] (0,2.6) to[R,l={#2}] (3,2.6) to[short,*-] (6,2.6) to[R,l={#4}] (6,0)
       to[short,-*] (3,0) to[short] (0,0);
 \draw (3,2.6) to[R,l={#3}] (3,0);
\end{circuitikz}}

% Mixt C: (R1 en sèrie amb R2) en paral·lel amb R3
\newcommand{\cmixc}[4]{%
\begin{circuitikz}[scale=0.85,transform shape]
 \draw (0,0) to[battery1,l_={#1}] (0,2.6) to[short] (2,2.6) to[R,l={#4}] (2,0) to[short] (0,0);
 \draw (2,2.6) to[short,*-] (4,2.6) to[R,l={#2}] (7,2.6) to[R,l={#3}] (10,2.6) to[short] (10,0) to[short] (2,0);
 \node[circ] at (2,0) {};
\end{circuitikz}}

% Mixt D: R1 en sèrie amb (R2 en paral·lel amb R3) i amb R4
\newcommand{\cmixd}[5]{%
\begin{circuitikz}[scale=0.8,transform shape]
 \draw (0,0) to[battery1,l_={#1}] (0,2.6) to[R,l={#2}] (3,2.6) to[short,*-] (6,2.6) to[R,l={#4}] (6,0)
       to[short,-*] (3,0) to[R,l_={#5}] (0,0);
 \draw (3,2.6) to[R,l={#3}] (3,0);
\end{circuitikz}}
